%% theory/eigen/diameter.tex -- Task 1: an explicit diameter bound on C(A, eps, M).
%% LaTeX-ready fragment in the notation of paper/jga/manuscript.tex (macros \Orb, \Area, \diam, ...).
%% Verification: theory/eigen/diameter.py (every displayed identity of hyperbolic trigonometry is
%% checked there on random configurations at 40 digits; Lemma eig:sep and Theorem eig:diam are
%% compared with enumerated triangle groups; Proposition eig:233 is checked for 7 <= m <= 60).

\subsection{The diameter is controlled by area, systole and the largest cone order}\label{sec:eig-diam}

For $A>0$, $\varepsilon>0$ and an integer $M\ge2$ let $\Cl(A,\varepsilon,M)$ be the class of $\Orb=\Gamma\backslash\mathbb H^2\in\Sig$ with $\Area(\Orb)\le A$, systole $\ell(\Orb)\ge\varepsilon$ (the least translation length of a hyperbolic element of $\Gamma$) and all cone orders at most $M$. A point $\tilde x\in\mathbb H^2$ with nontrivial stabiliser $\Gamma_{\tilde x}$ is an \emph{elliptic point}; $\Gamma_{\tilde x}$ is cyclic of order $m\le M$, generated by the rotation by $2\pi/m$ about $\tilde x$, and every elliptic element of $\Gamma$ is a nontrivial rotation about an elliptic point. We use the following facts of plane hyperbolic geometry (checked numerically in \texttt{diameter.py}):
\begin{enumerate}[label=(H\arabic*)]
\item a rotation by $\phi$ about $q$ moves a point at distance $r$ from $q$ by $d$, where $\sinh\frac d2=\sinh r\,|\sin\frac\phi2|$; a hyperbolic $\gamma$ moves every point by at least $\ell(\gamma)$;
\item the product of the reflections in two lines is a rotation by twice their angle if they meet, parabolic if they are asymptotic, and a translation by twice their distance if they are ultraparallel;
\item let $\lambda$ be the line through $p\ne q$, $d=d(p,q)$, $0<\alpha,\beta\le\frac\pi2$, and $L_1\ni p$, $L_3\ni q$ the lines making angles $\alpha$, $\beta$ with $\lambda$ on the same side of $\lambda$; put $X=\sin\alpha\sin\beta\cosh d-\cos\alpha\cos\beta$. If $X<1$, $L_1$ and $L_3$ meet at an angle $\theta\in(0,\pi)$ with $\cos\theta=X$, forming a triangle with angles $\alpha,\beta,\theta$ (the law of cosines for angles); if $X=1$ they are asymptotic; if $X>1$ they are ultraparallel at distance $h$ with $\cosh h=X$.
\end{enumerate}

\begin{lemma}[Cone points are separated]\label{eig:sep}
Let $\Orb\in\Cl(A,\varepsilon,M)$. Any two distinct elliptic points $\tilde p\ne\tilde q$ of $\Gamma$ satisfy
\[
  d(\tilde p,\tilde q)\ \ge\ d_0(\varepsilon,M):=\min\Big\{\frac\varepsilon2,\ \operatorname{arccosh}\Big(1+\frac2{\pi^2M^2}\Big)\Big\}.
\]
\end{lemma}

\begin{proof}
Let $a,b\in[2,M]$ be the orders of $\Gamma_{\tilde p}$, $\Gamma_{\tilde q}$, $\alpha=\pi/a$, $\beta=\pi/b$, $d=d(\tilde p,\tilde q)$, and $\lambda,L_1,L_3,X$ as in (H3), with reflections $s_\lambda,s_1,s_3$. By (H2), $s_1s_\lambda$ and $s_\lambda s_3$ are rotations about $\tilde p$, $\tilde q$ by $2\alpha$, $2\beta$; both senses of rotation lie in the cyclic stabilisers, so after replacing them by their inverses if necessary they lie in $\Gamma$, and so does $\gamma=s_1s_3$. As $\Gamma$ is cocompact it has no parabolic element, so $X\ne1$.

If $X>1$, $\gamma$ is hyperbolic with $\ell(\gamma)=2h$, $\cosh h=X\le\cosh d$ since $\cos\alpha\cos\beta\ge0$; so $\varepsilon\le\ell(\gamma)\le2d$.

If $X<1$, $L_1,L_3$ meet at $w$ at the angle $\theta$, with $\alpha+\beta+\theta<\pi$, and $\gamma$ is a nontrivial rotation about $w$ by $2\theta$. So $w$ is elliptic with stabiliser of order $c\le M$ containing a rotation by $2\theta$, whence $\theta=\pi k/c$, $1\le k<c$. By (H3),
\[
  \cosh d-1=\frac{\cos\theta+\cos(\alpha+\beta)}{\sin\alpha\sin\beta}=\frac{2\cos\frac{\theta+\alpha+\beta}2\,\cos\frac{\theta-\alpha-\beta}2}{\sin\alpha\sin\beta}.
\]
The defect is $\pi-\alpha-\beta-\theta=\pi\,(abc-bc-ac-kab)/(abc)\ge\pi/(abc)$, its numerator being a positive integer, so $\cos\frac{\theta+\alpha+\beta}2=\sin\frac{\pi-\alpha-\beta-\theta}2\ge\frac1{abc}$ by $\sin x\ge\frac{2x}\pi$ on $[0,\frac\pi2]$. Next $\theta\ge\pi/c\ge\pi/M$ and $\pi/M\le\alpha,\beta\le\pi/2$: if $\theta\ge\alpha+\beta$ then $0\le\theta-\alpha-\beta<\pi-2(\alpha+\beta)\le\pi-\frac{4\pi}M$, and otherwise $0<\alpha+\beta-\theta\le\pi-\frac\pi M$; so $\cos\frac{\theta-\alpha-\beta}2\ge\sin\frac\pi{2M}\ge\frac1M$. With $\sin\alpha\sin\beta\le\pi^2/(ab)$,
\[
  \cosh d-1\ \ge\ 2\cdot\frac1{abc}\cdot\frac1M\cdot\frac{ab}{\pi^2}=\frac2{\pi^2cM}\ \ge\ \frac2{\pi^2M^2}.\qedhere
\]
\end{proof}

\begin{lemma}[Balls have area bounded below]\label{eig:balls}
Let $\Orb\in\Cl(A,\varepsilon,M)$, $r_0=\frac12d_0(\varepsilon,M)$ and $\rho_1=\operatorname{arsinh}(\sinh r_0\sin\frac\pi M)$. Every metric ball of radius $2r_0$ in $\Orb$ has area at least
\[
  v_0(\varepsilon,M):=2\pi\min\Big\{\frac{\cosh r_0-1}M,\ \cosh\rho_1-1\Big\}.
\]
\end{lemma}

\begin{proof}
Let $\pi:\mathbb H^2\to\Orb$ be the projection, so $B_\Orb(\pi\tilde y,r)=\pi(B(\tilde y,r))$. If $d(\tilde y,\gamma\tilde y)\ge2r$ for every $\gamma\notin\Gamma_{\tilde y}$, then $\pi$ induces an injection of $\Gamma_{\tilde y}\backslash B(\tilde y,r)$ onto $B_\Orb(\pi\tilde y,r)$ (if $\gamma z=z'$ with $z,z'\in B(\tilde y,r)$ then $d(\tilde y,\gamma\tilde y)<2r$), and the ball has area $2\pi(\cosh r-1)/|\Gamma_{\tilde y}|$. Let $x\in\Orb$ and $\Sigma$ the set of cone points.

If $d(x,\Sigma)<r_0$, take $p\in\Sigma$ with $d(x,p)<r_0$ and a lift $\tilde p$. For $\gamma\notin\Gamma_{\tilde p}$, $\gamma\tilde p\ne\tilde p$ is elliptic, so $d(\tilde p,\gamma\tilde p)\ge d_0=2r_0$ (Lemma~\ref{eig:sep}); hence $B_\Orb(x,2r_0)\supset B_\Orb(p,r_0)$ has area at least $2\pi(\cosh r_0-1)/M$.

If $d(x,\Sigma)\ge r_0$, then $\Gamma_{\tilde x}=1$. A hyperbolic $\gamma$ moves $\tilde x$ by at least $\varepsilon\ge2d_0\ge2\rho_1$. An elliptic $\gamma$ is a rotation by $2\pi j/k$, $1\le j<k\le M$, about some $\tilde q$ with $d(\tilde x,\tilde q)\ge r_0$, and moves $\tilde x$ by $2\operatorname{arsinh}(\sinh d(\tilde x,\tilde q)\sin\frac{\pi j}k)\ge2\rho_1$ by (H1), as $\sin\frac{\pi j}k\ge\sin\frac\pi k\ge\sin\frac\pi M$. So $B_\Orb(x,\rho_1)$ has area $2\pi(\cosh\rho_1-1)$, and $\rho_1\le r_0$.
\end{proof}

\begin{theorem}[Diameter bound]\label{eig:diam}
Every $\Orb\in\Cl(A,\varepsilon,M)$ satisfies
\[
  \diam(\Orb)\ <\ D(A,\varepsilon,M):=\frac{4r_0A}{v_0(\varepsilon,M)}\ \le\ \frac A\pi\,\max(4M,M^2)\,\max\Big(\frac4\varepsilon,\frac{10M}3\Big).
\]
\end{theorem}

\begin{proof}
Let $\Delta=\diam\Orb=d(x,y)$ and choose lifts with $d(\tilde x,\tilde y)=\Delta$. The projection $c$ of the segment $[\tilde x,\tilde y]$, parametrised by arclength, satisfies $d(c(s),c(s'))=|s-s'|$: ``$\le$'' as $\pi$ is $1$-Lipschitz, and ``$\ge$'' since otherwise $d(x,y)\le s+d(c(s),c(s'))+(\Delta-s')<\Delta$. The points $c(4r_0i)$, $0\le i\le k:=\lfloor\Delta/4r_0\rfloor$, are $4r_0$ apart, so the balls of radius $2r_0$ about them are disjoint, and Lemma~\ref{eig:balls} gives $(k+1)v_0\le\Area(\Orb)\le A$; as $\Delta/4r_0<k+1$, $\Delta<4r_0A/v_0$.
For the closed form: $\cosh r-1\ge r^2/2$; $\cosh\rho_1-1=\sinh^2\rho_1/(\cosh\rho_1+1)\ge(\cosh r_0-1)\sin^2\frac\pi M$ and $\sin\frac\pi M\ge\frac2M$, so $v_0\ge\pi r_0^2/\max(M,M^2/4)$; and $\operatorname{arccosh}(1+x)\ge\sqrt{2x/(1+x)}$ (from $\cosh y-1\le\frac{y^2}2\cosh y$) gives $d_0\ge\min(\frac\varepsilon2,\frac{0.6}M)$.
\end{proof}

Nothing here uses a collar or thick--thin decomposition: the systole bounds the injectivity radius away from cone points, and Lemma~\ref{eig:sep} does so near them. The bound is crude (Table~\ref{eig:tab:D}); it grows like $A/\varepsilon$ as $\varepsilon\to0$ and like $AM^3$ as $M\to\infty$.

\begin{table}[h]
\centering\footnotesize
\caption{$D(A,\varepsilon,M)$ of Theorem~\ref{eig:diam}, and the diameters it bounds: $\Orb(2,8,8)$ and $\Orb(3,3,12)$ ($A=\pi/2$) have $\diam<4.9$ and $<4.4$; the eight members of the family of Figure~\ref{fig:F6} ($A=4\pi/3$) have $\diam\le7.8$.}\label{eig:tab:D}
\begin{tabular}{@{}lllll@{}}
\toprule
$A$ & $\varepsilon$ & $M=3$ & $M=12$ & $M=100$\\
\midrule
$\pi/2$ & $1$ & $56.6$ & $1125$ & $6.37\times10^5$\\
$4\pi/3$ & $1$ & $150.9$ & $3001$ & $1.70\times10^6$\\
$4\pi/3$ & $0.1$ & $640$ & $3184$ & $1.70\times10^6$\\
$10\pi$ & $1$ & $1132$ & $2.25\times10^4$ & $1.27\times10^7$\\
\bottomrule
\end{tabular}
\end{table}

\paragraph{The dependence on $M$ is necessary.} Near a cone point of large order the orbifold looks like a cusp, and its diameter grows.

\begin{proposition}[The family $\Orb(2,3,m)$]\label{eig:233}
For $m\ge7$ put $s_m=\operatorname{arccosh}\frac{2\cos(\pi/m)}{\sqrt3}$, $s_\infty=\operatorname{arccosh}\frac2{\sqrt3}$, $h_m=\operatorname{arccosh}\frac1{2\sin(\pi/m)}$ and
\[
  \sigma_0=2\operatorname{arsinh}\frac1{2\sqrt{1+\frac34\cosh^2(2s_\infty)}}=0.56206\ldots
\]
Then $\Area\Orb(2,3,m)=2\pi(\frac16-\frac1m)<\frac\pi3$, $\ell(\Orb(2,3,m))\ge\sigma_0$, $\diam\Orb(2,3,m)\ge h_m\ge\log\frac m{2\pi}$, and the ball of radius $h_m$ about the cone point of order $m$ is a hyperbolic cone of angle $2\pi/m$. Hence $\Orb(2,3,m)\in\Cl(\frac\pi3,\sigma_0,m)$ for all $m\ge7$, and any bound $\diam\le D(A,\varepsilon,M)$ valid on these classes has $D(A,\varepsilon,M)\ge\log\frac M{2\pi}$ for $A\ge\frac\pi3$, $\varepsilon\le\sigma_0$, $M\ge7$.
\end{proposition}

\begin{proof}
$\Orb=\Orb(2,3,m)$ is the double of the triangle $T$ with vertices $v_2,v_3,v_m$ and angles $\frac\pi2,\frac\pi3,\frac\pi m$; $\Gamma_m$ is the orientation-preserving subgroup of the group $\Gamma^*$ generated by the reflections in the sides of $T$, and the folding $f:\mathbb H^2\to T$ ($f(z)\in\Gamma^*z\cap T$) is $1$-Lipschitz with $d(z,\Gamma^*v)=d(f(z),v)$ for each vertex $v$, as $f(z)=gz$ for some $g\in\Gamma^*$. By the law of cosines for angles, $d(v_2,v_3)=s_m<s_\infty$ and $d(v_2,v_m)=h_m$. The folding of $\Orb$ onto $T$ is the identity on $T$ and $1$-Lipschitz, so the cone points of orders $2$ and $m$ are at distance $h_m$ in $\Orb$; and $h_m\ge\log\frac1{2\sin(\pi/m)}\ge\log\frac m{2\pi}$.

\emph{The cone ball.} Let $P$ be the union of the $2m$ tiles $gT$, $g\in\Gamma^*$, with vertex $\tilde v_m=v_m$: the union of the $\Gamma_{m,v_m}$-translates of a fundamental domain $T\cup s(T)$, $s$ the reflection in $[v_m,v_2]$. For $\gamma\in\Gamma_m\setminus\Gamma_{m,v_m}$, $\gamma P$ and $P$ have disjoint interiors. $P$ is star-shaped about $v_m$ and its boundary consists of images of $[v_2,v_3]$, whose distance from $v_m$ is $h_m$, because the angle of $T$ at $v_2$ is $\frac\pi2$. Hence $B(v_m,h_m)\subset\operatorname{int}P$, its translates are equal or disjoint, and its image is the cone of angle $2\pi/m$ and radius $h_m$.

\emph{The systole.} Let $\gamma\in\Gamma_m$ be hyperbolic with axis $\alpha$ and $L=\ell(\gamma)$. As $\alpha$ is connected and unbounded, it is not covered by the disjoint open balls $gB(v_m,h_m)$; pick $\tilde x\in\alpha$ outside all of them, and put $x=f(\tilde x)$, so $d(x,v_m)\ge h_m$. Write $x\in[v_m,q]$ with $q\in[v_2,v_3]$. Since $d(v_m,\cdot)$ is convex on $[v_2,v_3]$ with minimum $h_m=d(v_m,v_2)$ at $v_2$, $d(x,q)=d(v_m,q)-d(v_m,x)\le d(v_m,v_3)-d(v_m,v_2)\le s_m$, so $d(x,v_3)\le2s_m$. Hence some elliptic point $\tilde w\in\Gamma^*v_3=\Gamma_mv_3$, fixed by a rotation $\beta\in\Gamma_m$ by $\frac{2\pi}3$, lies at distance $r\le2s_m<2s_\infty$ from $\alpha$. The group $\langle\gamma,\beta\rangle$ is discrete; it is non-elementary, since a finite orbit in $\mathbb H^2\cup\partial\mathbb H^2$ would be a nonempty $\gamma$-invariant subset of the two endpoints of $\alpha$, which $\beta$, a rotation of order $3$, neither fixes nor swaps. Jørgensen's inequality \cite{jorgensen1976} gives $|\operatorname{tr}^2\gamma-4|+|\operatorname{tr}[\gamma,\beta]-2|\ge1$, where (in $SL(2,\R)$, signs immaterial)
\[
  \operatorname{tr}^2\gamma-4=4\sinh^2\tfrac L2,\qquad \operatorname{tr}[\gamma,\beta]-2=4\sinh^2\tfrac L2\,\sin^2\tfrac\pi3\,\cosh^2r
\]
(the second identity, for a rotation by $\phi$ at distance $r$ from the axis, with $\sin^2\frac\phi2$ in place of $\sin^2\frac\pi3$, is checked in \texttt{diameter.py}). So $4\sinh^2\frac L2\,(1+\frac34\cosh^2(2s_\infty))\ge1$, that is, $L\ge\sigma_0$.
\end{proof}

So the diameter must grow at least logarithmically in $M$; Theorem~\ref{eig:diam} gives $O(M^3)$. We do not determine the true order of growth.
